\subsection{Keystroke Dynamics in HCI}

Keystroke dynamics refers to the analysis of typing patterns --- not just what is typed, but how it is typed. Rather than treating typing as a uniform output channel, keystroke dynamics captures the rhythm, timing, and variability of each keypress, including inter-key intervals, hold times, burst patterns, and correction behaviour \parencite{epp2011identifying}. These micro-level temporal features have been shown to carry information about the typist's cognitive, emotional, and physiological state.

Early work established that keystroke latencies can reliably distinguish between individuals, leading to applications in biometric authentication \parencite{joyce1990identity}. Subsequent research extended this paradigm beyond identity verification to the detection of transient internal states. For example, keystroke features such as pause frequency, typing speed, and error rate have been used to identify emotional states including frustration, excitement, and boredom \parencite{epp2011identifying}. Similarly, increased pause duration and decreased typing fluency have been linked to cognitive load and stress \parencite{vizer2009automated}.

A particularly rich line of research examines how writing difficulties manifest in keystroke patterns. Neurodivergent populations, including individuals with dyslexia, exhibit distinct typing signatures: they pause more frequently and for longer durations --- both within and between words --- compared to neurotypical controls, even when their raw typing speed is comparable \parencite{sumner2013children}. These pauses are not simply motor delays but reflect higher-level cognitive bottlenecks in spelling, planning, and transcription processes. The finding is significant because it demonstrates that keystroke dynamics can serve as a behavioural proxy for cognitive processes that are otherwise invisible to the observer.

For ADHD populations specifically, typing challenges extend beyond pausing. Difficulties with working memory mean that a thought can be lost mid-sentence, producing an abrupt pause followed by deletion or restructuring of the preceding text. Proprioceptive differences --- a reduced awareness of hand and finger positioning --- can lead to increased typo rates, repeated corrections of the same character, and erratic inter-key intervals. Perfectionism and rejection-sensitive dysphoria (RSD), both common in ADHD, can further distort typing patterns by triggering excessive backspacing, sentence-level rewrites, and prolonged pauses at clause boundaries as the writer re-evaluates word choice before committing to text.

From a methodological standpoint, these findings position keystroke dynamics as a non-invasive, continuous, and ecologically valid sensor for cognitive state. Unlike surveys or lab-based cognitive tests, keystroke data can be collected passively during natural computer use. This makes it particularly suited for adaptive intervention systems that respond to the user's real-time cognitive and emotional state without imposing additional interaction overhead.

In the context of the present work, keystroke dynamics provides the sensing layer for the adaptive writing support system. The features extracted --- inter-key intervals, pause durations, burst rates, correction frequencies, and rhythm consistency --- are drawn directly from the literature discussed above and serve as the input from which the system infers the user's current writing state.
